\begin{myexem}
  \label{exem:run3}
  We continue the running example; see Example~\ref{exem:run1} and Example~\ref{exem:run2}.

  For all $d$, $\pmu_{212} = \pmu_{323} = \pmu_{344} = \pmu_{431} = 0$
  hence the node 4 is ``unused'' by the dual.
  For $2d = 2, 4$, $\pmu_{123} = \pmu_{231} = 0$ so
  the node $2$ is ``unused'' for low degree.

  At first, one could think that the dual variables can be used to reduce the systems,
  e.g. remove nodes or edges.
  However, it would be a mistake to remove the node 2 since the periodic trajectory with highest growth rate uses this node.



  It is also interesting to notice that the matrices corresponding to the dual variables have low rank.
  For example, for $2d = 2$, $\pmu_{131}$
  (resp. $\pmu_{312}$, $\pmu_{331}$) is the Dirac measure $5.873 \cdot \delta_{(0.917,0.399)}$
  (resp. $3.966 \cdot \delta_{(0.875,0.485)}$, $6.704 \cdot \delta_{(0.757,-0.653)}$).
  However, this is not a feasible solution of \progref{dualinf}.
  Indeed,
  while \eqref{eq:dualinf1} is satisfied for \node{} 1 since
  $\pushf{A_2}{\delta_{(0.875,0.485)}}$ gives $\delta_{(0.917,0.399)}$,
  \eqref{eq:dualinf1} is not satisfied for \node{} 3 as
  $\pushf{A_1}{\delta_{(0.917,0.399)}}$ gives $\delta_{(0.999, -0.0271)}$ and
  $\pushf{A_1}{\delta_{(0.757,-0.653)}}$ gives $\delta_{(0.422, -0.906)}$.
\end{myexem}
